% ============================ CHAPTER 1 — INTRODUCTION ============================
% Length: 3-6 pages. Present ONLY the problem & objectives, NOT the solution detail.

% GUIDE (chapter opening): 1-2 sentences stating what this chapter covers.
This chapter introduces the context and motivation of the project, defines its
objectives and scope, outlines the tentative solution, and describes the
organization of the thesis.

\section{Motivation}
% GUIDE: real-world situation -> problem/task -> who benefits. NO solution here.
% For this project: small/medium distributors manage purchasing, warehousing,
% sales invoicing and last-mile delivery with disconnected Excel files and paper.
% Consequences: stock inaccuracy, overselling, expired perishable goods (no FEFO),
% no audit trail, slow approvals. Emphasise urgency/scale (e.g. FMCG distribution).
Small and medium-sized distributors play an important role in moving goods from
manufacturers and wholesalers to retail stores, pharmacies, grocery shops, and
end customers. Their daily work includes purchasing goods from suppliers,
receiving stock into warehouses, monitoring inventory, selling to customers,
issuing invoices, and coordinating delivery. In many companies of this scale,
these activities are still handled by separated spreadsheets, paper forms,
phone messages, and manual approval. Such an approach can work when the business
is small, but it becomes fragile when the number of products, warehouses, sales
orders, and delivery trips increases.

The main problem is that operational data is fragmented. A purchase staff member
may record an incoming order in one file, the warehouse may record receipt in
another file, and the sales team may check stock by asking warehouse staff
directly. This creates delays and inconsistencies. Sales orders can be accepted
even when available stock is already committed to other customers. Goods with
expiry dates can remain in the warehouse because employees pick the most visible
stock instead of the earliest-expiring stock. When a delivery fails, the company
must manually restore stock, cancel or adjust invoices, and update order status.
These manual corrections are slow and error-prone.

The issue is especially serious for fast-moving consumer goods and perishable
products. In these domains, stock quantity alone is not enough; each batch must
be tracked with its expiry date, and outbound stock should follow a FEFO
principle. Without a reliable audit trail, the company has difficulty answering
basic questions such as which lot was delivered to which customer, why available
stock differs from physical stock, or whether an invoice still corresponds to a
successful delivery. Therefore, a unified information system is needed to support
daily distribution operations and reduce the risks caused by disconnected
manual management.

\section{Objectives and Scope of the Graduation Thesis}
% GUIDE: (1) overview existing products (Odoo, SAP Business One, KiotViet, Sapo,
% MISA), (2) compare/evaluate broadly, (3) state their limitations for a small
% distributor, (4) state what THIS thesis builds: an integrated web system covering
% purchasing, inventory with FEFO lot tracking, sales & invoicing, and delivery,
% with role-based access control. List the MAIN functions only (details -> Ch.5).
Several existing products already support inventory and distribution
management. Odoo provides inventory functions such as expiration-date tracking
and FEFO removal strategy \cite{odooExpiration,odooFefo}. SAP Business One
supports batch master data with fields such as expiration date and manufacturing
date \cite{sapBatch}. Vietnamese retail-oriented systems such as KiotViet,
Sapo, and MISA also provide lot and expiry-date management features
\cite{kiotvietLotExpiry,sapoLotExpiry,misaLotExpiry}. These products show that
the business need is real and that lot-based inventory is an important function
for distribution companies.

However, such products also have limitations for the target context of this
thesis. Full ERP solutions are often costly, complex to customize, and excessive
for a small distributor that needs a focused workflow. Retail POS products are
convenient for store operations but may not fully match a distributor's combined
workflow of purchase approval, warehouse receiving, sales order approval,
inventory reservation, goods issue, delivery trip execution, and invoice
recovery after failed delivery. The objective of this graduation thesis is
therefore to design and implement a focused web-based supply-chain management
system for small and medium-sized distributors.

The scope of the thesis includes five main functional groups. The purchasing
module supports supplier management, purchase order creation, approval, and
goods receipt. The inventory module supports multi-warehouse stock records,
inventory transactions, lot tracking, expiry dates, and FEFO-based outbound
allocation. The sales module supports customer management, sales order creation,
approval with stock reservation, goods issue, and sales invoice generation. The
delivery module supports delivery plans, delivery trips, shipper assignment, and
delivery status updates. The administration module supports authentication,
role-based access control, user management, and dashboard reporting. Advanced
functions such as automatic route optimization, demand forecasting, barcode
scanning, and a dedicated mobile application are considered future work.

\section{Tentative Solution}
% GUIDE: (i) name the approach/tech (layered Spring Boot REST backend + React SPA,
% PostgreSQL, JWT/RBAC, FEFO lot allocation); (ii) 1-2 sentence description;
% (iii) main contribution & expected results. Do NOT analyse the tech in detail.
The proposed solution is a web application based on a layered client-server
architecture. The backend is implemented with Spring Boot REST services and
PostgreSQL, while the frontend is implemented as a React single-page
application. Authentication and authorization are handled by JWT and role-based
access control so that each actor can access only the functions required by his
or her responsibility.

The core business direction is to connect procure-to-stock, order-to-cash, and
delivery execution in one consistent workflow. Goods receipt creates inventory
lots with batch numbers and expiry dates. Sales order approval reserves
available stock to prevent overselling. Goods issue confirms the physical
outbound movement and allocates stock by FEFO. When delivery fails, the system
restores inventory and adjusts the related order and invoice state. The expected
result is a working prototype that demonstrates how a small distributor can
control stock, approvals, deliveries, and traceability in an integrated system.

\section{Thesis Organization}
% GUIDE: full paragraphs (not bullets); describe Chapters 2-6, NOT Chapter 1.
Chapter~\ref{chapter:Survey} surveys existing distribution systems and analyses the
functional and non-functional requirements. Chapter~\ref{chapter:Background}
presents the theoretical background and technologies adopted.
Chapter~\ref{chapter:Design} describes the architecture, detailed design,
implementation and evaluation of the system. Chapter~\ref{chapter:Contribution}
discusses the main contributions, in particular the FEFO inventory-allocation
mechanism. Chapter~\ref{chapter:Conclusion} concludes the thesis and outlines
future work.
